\documentclass[12pt,a4paper]{article}

\usepackage[a4paper,text={16.5cm,25.2cm},centering]{geometry}
\usepackage{lmodern}
\usepackage{amssymb,amsmath}
\usepackage{bm}
\usepackage{graphicx}
\usepackage{microtype}
\usepackage{hyperref}
\usepackage{minted}
\setlength{\parindent}{0pt}
\setlength{\parskip}{1.2ex}

\hypersetup
       {   pdfauthor = {  },
           pdftitle={  },
           colorlinks=TRUE,
           linkcolor=black,
           citecolor=blue,
           urlcolor=blue
       }






\begin{document}



\section{Inverzna potenčna metoda za zgornjo Hessenbergovo matriko}
Tina Brdnik

\subsection{Opis naloge}
V poročilu je opisana implementacija inverzne potenčne metode za določanje lastnih vektorjev matrike.  Metoda temelji na zaporedju: 

\[
x^{(n+1)} = \frac{A_{\lambda}^{-1} \, x^{(n)}}{\left\| A_{\lambda}^{-1} x^{(n)} \right\|_{\infty}}, \quad A_{\lambda} = A - \lambda I
\]
ki konvergira k lastnemu vektorju, ki pripada lastni vrednosti matrike $A$ najbližji \ensuremath{\lambda}.

Za zmanjšanje števila operacij, matriko $A$ najprej pretvorimo v zgornjo  Hessenbergovo obliko, s Householderjevimi zrcaljenji, saj ta ohrani lastne  vrednosti in hkrati poenostavi LU razcep, saj ima matrika redko obliko.

\subsection{Opis rešitve}
Implementacija algoritma poteka v več korakih:

\begin{itemize}
\item[1. ] Pretvorba matrike v Hessenbergovo obliko  

\end{itemize}
Matriko $A$ najprej pretvorimo v zgornjo Hessenbergovo obliko s pomočjo Householderjevih zrcaljenj. S tem odpravimo elemente pod drugo poddiagonalo in dobimo matriko $H$, za katero velja $H = Q^T A Q$, pri čemer je $Q$ ortogonalna matrika. Ker imata $A$ in $H$ enake lastne vrednosti, lahko nadaljnji izračun izvajamo na $H$.  

Prednost Hessenbergove oblike je, da omogoča učinkovitejši LU razcep, saj je matrika $L$ spodnje tridiagonalna, matrika $U$ pa zgornjetrikotna. Tako se časovna zahtevnost razcepa zmanjša iz $\mathcal{O}(n^3)$ na $\mathcal{O}(n^2)$.

\begin{itemize}
\item[2. ] LU razcep  

\end{itemize}
Za reševanje sistemov izvedemo LU razcep zgornje Hessenbergove matrike. Pri tem ostane $U$ zgornjetrikotna, $L$ pa spodnje tridiagonalna z enicami na diagonali. Zaradi redke strukture poteka eliminacija le ob diagonali, kar zagotavlja časovno zahtevnost $\mathcal{O}(n^2)$.

\begin{itemize}
\item[3. ] Inverzna iteracija  

\end{itemize}
Inverzna potenčna metoda poišče lastno vrednost matrike, ki je najbližja podanemu približku. Na vsakem koraku rešimo sistem $(A - \lambda I)y = x$ in dobljeni vektor normiramo. Ker račun poteka na Hessenbergovi obliki in uporabljamo LU razcep, je metoda bistveno hitrejša kot v splošnem primeru.

\begin{itemize}
\item[4. ] Rekonstrukcija lastnega vektorja  

\end{itemize}
Če je $v$ lastni vektor Hessenbergove matrike $H$, potem je $Qv$ lastni vektor izvorne matrike $A$.

\subsection{Primer uporabe}
Inverzno iteracijo smo preizkusili na matriki polinoma:

\[
p(x) = x^3 - 2x^2 - 4x + 6,
\]
katere lastne vrednosti ustrezajo ničlam polinoma.   Na spodnji sliki je narisan graf polinoma, ničle in poti konvergence za različne začetne približke $x_0$.


\begin{minted}[texcomments = true, mathescape, fontsize=\small, xleftmargin=0.5em]{julia}
include("demo.jl")
using .Naloga1
nariši_polinom_s_konvergenco([1.0, -2.0, -4.0, 6.0])
\end{minted}
\includegraphics[width=\linewidth]{jl_LWxXfN/porocilo_0_1.pdf}


\end{document}
